\documentclass[11pt,reqno]{amsart}
\usepackage[utf8]{inputenc}
\usepackage[T1]{fontenc}
\usepackage{amsmath,amssymb,amsthm}
\usepackage[margin=1in]{geometry}
\usepackage[colorlinks=true,linkcolor=blue,citecolor=blue,urlcolor=blue]{hyperref}

\newcommand{\F}{\mathbb F}
\newcommand{\Z}{\mathbb Z}
\newcommand{\G}{\mathfrak G}
\newcommand{\LL}{\mathfrak L}
\newcommand{\D}{\mathfrak D}
\newcommand{\Sm}{\mathcal S}
\newcommand{\Lh}{\widehat{\LL}}
\newcommand{\rp}[1]{\|#1\|}
\DeclareMathOperator{\spn}{span}

\theoremstyle{plain}
\newtheorem{theorem}{Theorem}
\newtheorem{lemma}[theorem]{Lemma}
\newtheorem{proposition}[theorem]{Proposition}
\newtheorem{corollary}[theorem]{Corollary}
\theoremstyle{definition}
\newtheorem{definition}[theorem]{Definition}
\theoremstyle{remark}
\newtheorem{remark}[theorem]{Remark}

\sloppy

\begin{document}

\title{A basis of free generalized Gerstenhaber algebras}

\author{Ivan Kaygorodov}
\address{CMA-UBI, Universidade da Beira Interior, Covilh\~a, Portugal}
\email{kaygorodov.ivan@gmail.com}

\author{Kseniia Viatkina}
\address{Saint Petersburg University, Saint Petersburg, Russia}
\email{vjatkina.k@gmail.com}

\thanks{The work is supported by RSF 22-71-10001.}

\subjclass[2020]{17B63, 17B01, 17A50}
\keywords{Gerstenhaber algebra, odd Poisson superalgebra, generalized Poisson superalgebra, free algebra, basis, Lie superalgebra}

\begin{abstract}
Generalized Gerstenhaber algebras are the odd generalized Poisson superalgebras of Cantarini and Kac; Gerstenhaber algebras (odd Poisson superalgebras) form the special case in which the map $\D(a)=\{1,a\}$ vanishes.
We construct a linear basis of the free unital generalized Gerstenhaber algebra generated by a linearly ordered set of even and odd generators, following the method used earlier for free generalized Poisson superalgebras.
As an application we show that the multilinear part of the free generalized Gerstenhaber algebra on $n$ generators has dimension $n\cdot n!$.
\end{abstract}

\maketitle

\section{Introduction}

The notion of a Poisson bracket has its origin in the works of S.~D.~Poisson on celestial mechanics at the beginning of the XIX century.
Since then Poisson algebras have appeared in several areas of mathematics, such as Poisson manifolds~\cite{L1}, algebraic geometry~\cite{BLLM,GK04}, noncommutative geometry~\cite{ber08}, operads, quantization theory~\cite{Kon03}, quantum groups, classical and quantum mechanics.
The study of Poisson algebras has also led to a number of related algebraic structures: generic Poisson algebras, algebras of Jordan brackets and generalized Poisson algebras, Novikov--Poisson algebras, Malcev--Poisson--Jordan algebras, transposed Poisson algebras, $n$-ary Poisson algebras, Gerstenhaber algebras~\cite{K96}, and others.

A Gerstenhaber algebra (also called an odd Poisson superalgebra, or an anti-Poisson algebra) is a supercommutative associative superalgebra equipped with a bracket of degree $-1$ which is a Lie superalgebra bracket with respect to the reversed parity and acts on the associative product by derivations.
Such brackets arise in physics, most notably in the Batalin--Vilkovisky formalism~\cite{BV81,p1,p2}, and in geometry~\cite{p3}.
One of the most important purely algebraic sources of Gerstenhaber algebras is the Hochschild cohomology of an associative algebra~\cite{Ger63}; for recent results on Gerstenhaber brackets on Hochschild cohomology see~\cite{W21,V16,V19,L21}.
Subalgebras of the Gerstenhaber algebra of differential operators are described in~\cite{B06}.

Generalized Poisson superalgebras appeared in the study of Jordan superalgebras.
Kantor~\cite{Kan90,Kan92} associated to a Poisson superalgebra a Jordan superalgebra (the Kantor double), and King and McCrimmon~\cite{KM92} described the brackets whose Kantor double is a Jordan superalgebra; the resulting superalgebras of Jordan brackets are, up to a change of the bracket, exactly the generalized Poisson superalgebras (see~\cite{CK07,K17}; a generalized version of the Kantor double was studied in~\cite{K10}).
In a generalized Poisson superalgebra the Leibniz rule holds only up to a correction term involving the linear map $\D(a)=\{1,a\}$.
Simple linearly compact generalized Poisson superalgebras were classified by Cantarini and Kac~\cite{CK07}, who later introduced odd generalized Poisson superalgebras and classified the simple linearly compact ones~\cite{CK10}.
Odd generalized Poisson superalgebras with $\D=0$ are precisely Gerstenhaber algebras; accordingly, we call odd generalized Poisson superalgebras \emph{generalized Gerstenhaber algebras}.

Bases of free algebras were constructed for free Lie algebras~\cite{Hall50,shirshov,chi,Reu93}, free Lie superalgebras~\cite{S86,BMPZ,BKLM}, free Poisson superalgebras~\cite{She93}, free generic Poisson algebras~\cite{KSU}, free generalized Poisson superalgebras and free superalgebras of Jordan brackets~\cite{K17}, and so on.
The aim of the present note is to construct a basis of the free unital generalized Gerstenhaber algebra generated by a linearly ordered set of even and odd generators (Theorem~\ref{basis}) and to compute the dimension of its multilinear part (Theorem~\ref{dimen}).
We follow the method of Shirshov~\cite{shirshov} in the form used in~\cite{K17} for free generalized Poisson superalgebras: the candidate basis $U$ consists of supercommutative monomials in the elements of a basis of the free odd Lie superalgebra, and the linear independence of $U$ is proved by constructing a generalized Gerstenhaber algebra structure on the linear span of $U$.
There is, however, a new phenomenon in the odd case: the element $\{1,1\}$, which vanishes automatically in every generalized Poisson superalgebra, is a nonzero basis element of the free generalized Gerstenhaber algebra, so that the map $\D$ is no longer a derivation.
To handle this we define the bracket on the linear span of $U$ by an explicit closed formula (Section~\ref{sec:model}), for which skew-symmetry and the Leibniz rule are verified directly, and we derive the Jacobi identity from a general formula for the Jacobiator (Lemma~\ref{jacobiator}), which may be of independent interest.

\section{Definitions and notation}

Throughout the paper $\F$ is a field of characteristic zero.
All vector spaces and algebras are $\Z_2$-graded, and the elements appearing in formulas involving parities are assumed to be homogeneous.
For a homogeneous element $a$ we write $|a|\in\Z_2$ for its parity and
$$
\rp{a}=|a|+1\in\Z_2
$$
for its parity with respect to the reversed $\Z_2$-grading.

\begin{definition}\label{def}
A \emph{generalized Gerstenhaber algebra} is a $\Z_2$-graded vector space $V=V_0\oplus V_1$ with an element $1\in V_0$ and two bilinear operations $(a,b)\mapsto ab$ and $(a,b)\mapsto\{a,b\}$ such that for all homogeneous $a,b,c\in V$
\begin{align}
&|ab|=|a|+|b|,\qquad |\{a,b\}|=|a|+|b|+1;\label{id:deg}\\
&(ab)c=a(bc),\qquad ab=(-1)^{|a||b|}ba,\qquad 1a=a;\label{id:comm}\\
&\{a,b\}=-(-1)^{\rp a\rp b}\{b,a\};\label{id:skew}\\
&\{a,\{b,c\}\}=\{\{a,b\},c\}+(-1)^{\rp a\rp b}\{b,\{a,c\}\};\label{id:jacobi}\\
&\{a,bc\}=\{a,b\}c+(-1)^{\rp a|b|}b\{a,c\}+(-1)^{\rp a}\D(a)bc,\qquad\text{where } \D(a)=\{1,a\}.\label{id:leibniz}
\end{align}
A generalized Gerstenhaber algebra with $\D=0$ is called a \emph{Gerstenhaber algebra}.
\end{definition}

Thus the product is even, associative and supercommutative with unit $1$, the bracket is odd, and \eqref{id:skew}, \eqref{id:jacobi} are the odd skew-symmetry and the odd Jacobi identity, while \eqref{id:leibniz} is a generalized odd Leibniz rule.
Since $\rp a\rp b=(|a|-1)(|b|-1)$ in $\Z_2$, Definition~\ref{def} is exactly the definition of an odd generalized Poisson superalgebra given in~\cite[(0.14)]{CK10}, and Gerstenhaber algebras are exactly odd Poisson superalgebras.
The basic examples of generalized Gerstenhaber algebras with $\D\neq0$ are the superalgebras $PO(n,n+1)$ of~\cite{CK10}.

\begin{remark}[reversed parity]\label{rem:parity}
Let $\Pi V$ denote the space $V$ with the reversed $\Z_2$-grading, so that the parity of $a$ in $\Pi V$ is $\rp a$.
Identities \eqref{id:deg}, \eqref{id:skew} and \eqref{id:jacobi} say precisely that $(\Pi V,\{\,,\})$ is a Lie superalgebra.
We call a superalgebra with an odd bilinear operation satisfying \eqref{id:skew} and \eqref{id:jacobi} an \emph{odd Lie superalgebra}.
Reversing the parity establishes a bijection between odd Lie superalgebras and Lie superalgebras (cf.~\cite{Kac77} and~\cite[Remark~1.3]{CK10}); in particular, the free odd Lie superalgebra generated by a set $X$ of even and odd elements is the free Lie superalgebra generated by the same set with reversed parities.
\end{remark}

\begin{remark}[the map $\D$]\label{rem:D}
(a) Taking $a=1$ in \eqref{id:leibniz} we obtain
\begin{equation}\label{Dprod}
\D(bc)=\D(b)c+(-1)^{|b|}b\D(c)-\D(1)bc .
\end{equation}
Hence $\D$ is an odd derivation of $(V,\cdot)$ if and only if $\D(1)=\{1,1\}=0$; in general, $\D-\D(1)\,\mathrm{id}$ is an odd derivation.
In a generalized Poisson superalgebra (where the bracket is even) the element $\{1,1\}$ vanishes by skew-symmetry, but \eqref{id:skew} imposes no condition on $\{1,1\}$ because $\rp 1=1$.
We shall see (Corollary~\ref{cor:11}) that $\{1,1\}\neq0$ in the free generalized Gerstenhaber algebra.

(b) Using \eqref{id:comm} and $|\D(a)|=\rp a$, one checks that \eqref{id:leibniz} is equivalent to the following statement: for every $a\in V$ the map
$$
d_a\colon b\mapsto \{a,b\}+(-1)^{\rp a}\D(a)b
$$
is a superderivation of parity $\rp a$ of the superalgebra $(V,\cdot)$.
In particular $d_a(1)=0$, i.e. $\{a,1\}=-(-1)^{\rp a}\D(a)$, in accordance with \eqref{id:skew}.
\end{remark}

\section{The free generalized Gerstenhaber algebra and the set \texorpdfstring{$U$}{U}}\label{sec:free}

Let $X=\{1,x_1,x_2,\dots\}$ be a finite or countable set, linearly ordered by $1<x_1<x_2<\cdots$, in which $1$ is even and each $x_i$ has a prescribed parity $|x_i|\in\Z_2$.
Let $D_X$ be the free algebra with two bilinear operations $\cdot$ and $\{\,,\}$ generated by $X$; its basis consists of the nonassociative words in the alphabet $X$ built with these two operations.
The words are homogeneous, with $|1|=0$, $|uv|=|u|+|v|$ and $|\{u,v\}|=|u|+|v|+1$, and the \emph{degree} of a word is the number of letters in it.
Let $I$ be the ideal of $D_X$ generated by the elements
\begin{gather*}
ab-(-1)^{|a||b|}ba,\qquad (ab)c-a(bc),\qquad 1a-a,\qquad \{a,b\}+(-1)^{\rp a\rp b}\{b,a\},\\
\{a,\{b,c\}\}-\{\{a,b\},c\}-(-1)^{\rp a\rp b}\{b,\{a,c\}\},\\
\{a,bc\}-\{a,b\}c-(-1)^{\rp a|b|}b\{a,c\}-(-1)^{\rp a}\{1,a\}bc,
\end{gather*}
where $a,b,c$ run over the homogeneous elements of $D_X$.
Clearly, $\G=D_X/I$ is the free unital generalized Gerstenhaber algebra generated by $X$: every parity preserving map from $\{x_1,x_2,\dots\}$ to a generalized Gerstenhaber algebra $V$ extends uniquely to a homomorphism $\G\to V$ sending $1$ to $1$.

\subsection*{Good words}
Consider the nonassociative words in the alphabet $X$ with respect to the single operation $\{\,,\}$.
We define the set of \emph{good words} together with a linear order on it by induction on the degree.
The good words of degree $1$ are the elements of $X$, ordered as above.
Suppose that the good words of degree less than $n$ are defined and linearly ordered.
A word $w$ of degree $n$ is good whenever
\begin{enumerate}
\item[(1)] $w=\{u,v\}$, where $u$ and $v$ are good words with $u>v$;
\item[(2)] if $u=\{u_1,u_2\}$, then $u_2\le v$.
\end{enumerate}
Every good word of degree $n$ is greater than every good word of smaller degree; two good words $w=\{u_1,u_2\}$ and $w'=\{u_1',u_2'\}$ of degree $n$ are compared lexicographically: $w>w'$ if $u_1>u_1'$, or $u_1=u_1'$ and $u_2>u_2'$.
In particular, if $w=\{u,v\}$ is a good word, then $w>u>v$.
Thus the good words are the basic commutators of M.~Hall~\cite{Hall50}, and they form a Hall set in the sense of~\cite[Section~4.1]{Reu93}.

\subsection*{The set \texorpdfstring{$M$}{M}}
Put
$$
M=\{\,u,\ \{v,v\}\ \mid\ u,v \text{ are good words and } |v|=0\,\}.
$$
We extend the order to $M$ by the same rule: elements of larger degree are greater, and elements of the same degree are compared lexicographically by their pairs of components, a square $\{v,v\}$ being compared as the pair $(v,v)$.
This is a linear order on $M$ (a square $\{v,v\}$ is never a good word, since the components $u_1>u_2$ of a good word are distinct), and $1$ is the least element of $M$.
We enumerate the elements of $M$ as $e_1<e_2<\cdots$, so that $e_1=1$.

Let $\LL_X$ be the free algebra with one bilinear operation $\{\,,\}$ generated by $X$ (with the parities of words defined as above) and let $J$ be the ideal of $\LL_X$ generated by the elements
$\{a,b\}+(-1)^{\rp a\rp b}\{b,a\}$ and $\{a,\{b,c\}\}-\{\{a,b\},c\}-(-1)^{\rp a\rp b}\{b,\{a,c\}\}$, $a,b,c\in\LL_X$ homogeneous.
The quotient $\LL=\LL_X/J$ is the free odd Lie superalgebra generated by $X$.

\begin{lemma}\label{lem:M}
The set $M$ is a basis of $\LL$.
\end{lemma}

\begin{proof}
By Remark~\ref{rem:parity}, $\LL$ is the free Lie superalgebra generated by the set $X$ in which the parity of $x\in X$ is $\rp x$.
It is well known that the free Lie superalgebra over a field of characteristic zero generated by a linearly ordered set has a basis consisting of the elements of a Hall set (in particular, of the basic commutators defined above) and the squares $\{v,v\}$ of its odd elements $v$, see~\cite{S86,BMPZ,BKLM} (cf.~\cite[Section~3]{K17}).
An element $v$ of $\LL$ is odd with respect to the reversed parity if and only if $|v|=0$; hence this basis is exactly $M$.
\end{proof}

\subsection*{The set \texorpdfstring{$U$}{U}}
Let
$$
U=\bigl\{\,e_{i_1}^{k_1}e_{i_2}^{k_2}\cdots e_{i_r}^{k_r}\ \bigm|\ r\ge0,\ 1<i_1<i_2<\dots<i_r,\ k_j\ge1,\ \text{and } k_j=1 \text{ if } |e_{i_j}|=1\,\bigr\}
$$
be the set of supercommutative monomials in the elements of $M\setminus\{1\}$ (for $r=0$ the empty product is $1$).
For an element $u=e_{i_1}^{k_1}\cdots e_{i_r}^{k_r}$ of $U$ we call $\deg_e(u)=\sum_j k_j$ its \emph{$e$-degree}; the elements of $M$ are the elements of $U$ of $e$-degree $1$ together with $1$.
We regard the elements of $M$ and $U$ as elements of $\G$ (evaluating the words in $\G$).
Since the product of $\G$ is supercommutative and $ee=0$ for odd $e$, the product of two elements of $U$ is, up to sign, again an element of $U$ or zero.

The main goal of the paper is to prove that $U$ is a basis of $\G$.
We begin with the easier half.

\begin{lemma}\label{vsp}
The set $U$ spans $\G$.
\end{lemma}

\begin{proof}
Since $(\G,\{\,,\})$ satisfies \eqref{id:skew} and \eqref{id:jacobi}, the identity map on $X$ extends to a homomorphism $\psi\colon\LL\to\G$ of odd Lie superalgebras; it sends each element of $M$ to the same element of $\G$.
By Lemma~\ref{lem:M}, $\psi(\LL)$ is spanned by $M\subset U$.

We show by induction on the degree $n$ of a word $w\in D_X$ that the image of $w$ in $\G$ lies in $\spn U$; this proves the lemma, since the words span $D_X$.
For $n=1$, $w\in X\subset M$.
Let $n\ge2$ and assume the claim for words of degree less than $n$.
Then $w=w_1w_2$ or $w=\{w_1,w_2\}$, where $w_1,w_2$ are words of degree less than $n$, so $w_1$ and $w_2$ lie in $\spn U$.
If $w=w_1w_2$, then $w\in\spn U$ because $U$ is closed under the product up to sign.
Let $w=\{w_1,w_2\}$; it suffices to show that $\{u,v\}\in\spn U$ for all $u,v\in U$ occurring in the expansions of $w_1$ and $w_2$, respectively.
If $v=v_1v_2$ with $v_1,v_2\in U\setminus\{1\}$, then by \eqref{id:leibniz}
$$
\{u,v_1v_2\}=\{u,v_1\}v_2+(-1)^{\rp u|v_1|}v_1\{u,v_2\}+(-1)^{\rp u}\{1,u\}v_1v_2 .
$$
Each of $\{u,v_1\}$, $\{u,v_2\}$ and $\{1,u\}$ is a linear combination of words of degree less than $n$, hence lies in $\spn U$ by the induction hypothesis, and $\spn U$ is closed under the product; therefore $\{u,v\}\in\spn U$.
If $u=u_1u_2$ with $u_1,u_2\in U\setminus\{1\}$, we first apply \eqref{id:skew} and then argue as above.
In the remaining case $u,v\in M$, so $\{u,v\}=\psi(\{u,v\})\in\psi(\LL)\subset\spn U$.
\end{proof}

\section{The model algebra}\label{sec:model}

In this section we construct a generalized Gerstenhaber algebra structure on the free supercommutative superalgebra generated by $M\setminus\{1\}$.

Let $\LL'=\spn(M\setminus\{1\})\subset\LL$, and let $\Sm$ be the free supercommutative superalgebra generated by the superspace $\LL'$, i.e. the tensor product of the symmetric algebra of $\LL'_0$ and the exterior algebra of $\LL'_1$.
The monomials in the elements of the basis $M\setminus\{1\}$ of $\LL'$, that is, the elements of the set $U$, form a basis of $\Sm$.
Put $\Lh=\F1\oplus\LL'\subset\Sm$.
The linear map $\LL\to\Lh$ which is the identity on $\LL'$ and sends the basis element $1\in M$ to the unit $1\in\Sm$ is an isomorphism of superspaces (by Lemma~\ref{lem:M}); we use it to transfer the bracket of $\LL$ to $\Lh$.
Thus $\Lh$ is an odd Lie superalgebra with basis $M$, in which $1$ is the least basis element, and $\Lh$ generates $\Sm$ as a unital algebra.
For $u\in\Lh$ we put $\D(u)=\{1,u\}\in\Lh$; note that $\D(1)=\{1,1\}$ is an odd element of $M\setminus\{1\}$.

\subsection*{Notation}
For a product $a=a_1a_2\cdots a_m$ of homogeneous elements of $\Sm$ and $1\le i\le m$ we write
$$
a_{<i}=a_1\cdots a_{i-1},\qquad a_{>i}=a_{i+1}\cdots a_m,\qquad a_{\hat\imath}=a_{<i}a_{>i}
$$
(empty products are equal to $1$).
By \eqref{id:comm}, $a_{<i}a_i=(-1)^{|a_i||a_{<i}|}a_ia_{<i}$ and $a=(-1)^{|a_i||a_{<i}|}a_ia_{\hat\imath}=(-1)^{|a_i||a_{>i}|}a_{\hat\imath}a_i$.

\subsection*{The bracket}
For $a=a_1\cdots a_m$ and $b=b_1\cdots b_n$ with $m,n\ge1$ and homogeneous $a_i,b_j\in\Lh$ we put
\begin{equation}\label{Ddef}
\D(a)=\sum_{i=1}^m(-1)^{|a_{<i}|}\,a_{<i}\{1,a_i\}a_{>i}-(m-1)\{1,1\}\,a
\end{equation}
and
\begin{multline}\label{brdef}
\{a,b\}^*=\sum_{i=1}^m\sum_{j=1}^n(-1)^{|a_i||a_{>i}|+|b_j||b_{<j}|}\,a_{\hat\imath}\{a_i,b_j\}b_{\hat\jmath}\\
-(m-1)\,a\D(b)+(n-1)(-1)^{\rp a}\D(a)b-(m-1)(n-1)\,a\D(1)b .
\end{multline}
Here $\{a_i,b_j\}$ is the bracket of $\Lh$, so that $\{a,b\}^*$ is defined in terms of the bracket of $\LL$ and the product of $\Sm$.
For $m=1$ and $a_1=1$ formula \eqref{Ddef} gives $\D(1)=\{1,1\}$, in agreement with the earlier notation, and formula \eqref{brdef} with $m=1$, $a_1=1$ gives $\{1,b\}^*=\D(b)$ (see Lemma~\ref{lem:welldef}).
The first sum in \eqref{brdef} is the bracket obtained by moving $a_i$ to the right end of $a$ and $b_j$ to the left end of $b$ and bracketing the two adjacent factors; the remaining terms are corrections involving $\D$.

\begin{lemma}\label{lem:welldef}
The right-hand sides of \eqref{Ddef} and \eqref{brdef} do not depend on the choice of the factorizations $a=a_1\cdots a_m$ and $b=b_1\cdots b_n$ with $a_i,b_j\in\Lh$.
Consequently, \eqref{Ddef} defines a linear map $\D\colon\Sm\to\Sm$ and \eqref{brdef} defines a bilinear map $\{\,,\}^*\colon\Sm\times\Sm\to\Sm$.
Moreover,
\begin{enumerate}
\item[(i)] $\{u,v\}^*=\{u,v\}$ and $\D(u)=\{1,u\}$ for $u,v\in\Lh$;
\item[(ii)] $\{1,a\}^*=\D(a)$ for all $a\in\Sm$;
\item[(iii)] $\D(bc)=\D(b)c+(-1)^{|b|}b\D(c)-\D(1)bc$ for all $b,c\in\Sm$;
\item[(iv)] $|\{a,b\}^*|=|a|+|b|+1$ and $|\D(a)|=\rp a$ for all homogeneous $a,b\in\Sm$.
\end{enumerate}
\end{lemma}

\begin{proof}
Fix $m$ and $n$.
The right-hand side of \eqref{Ddef} is a multilinear function $R_m(a_1,\dots,a_m)$ of $a_1,\dots,a_m\in\Lh$, and the right-hand side of \eqref{brdef} is a multilinear function $R_{m,n}(a_1,\dots,a_m;b_1,\dots,b_n)$.
We claim that these functions are
\begin{enumerate}
\item[(a)] supersymmetric: interchanging two neighbouring arguments $a_k,a_{k+1}$ (or $b_k,b_{k+1}$) multiplies the value by $(-1)^{|a_k||a_{k+1}|}$ (respectively, by $(-1)^{|b_k||b_{k+1}|}$);
\item[(b)] compatible with removing units: $R_m(a_1,\dots,a_{m-1},1)=R_{m-1}(a_1,\dots,a_{m-1})$, $R_{m,n}(a_1,\dots,a_{m-1},1;b)=R_{m-1,n}(a_1,\dots,a_{m-1};b)$ for $m\ge2$, and similarly in the second group of arguments.
\end{enumerate}
Let us verify (a) for $R_{m,n}$ and the arguments $a_k,a_{k+1}$; the other cases are analogous or simpler.
The products $a$, $a\D(b)$ and $a\D(1)b$ are multiplied by $(-1)^{|a_k||a_{k+1}|}$, and so is $\D(a)b$ once (a) is known for $R_m$.
In the double sum, the terms with $i\ne k,k+1$ are multiplied by $(-1)^{|a_k||a_{k+1}|}$, since $a_{\hat\imath}$ is and the exponent $|a_i||a_{>i}|$ is unchanged.
The term with $i=k$ of the original sum is $(-1)^{|a_k|(|a_{k+1}|+|a_{>k+1}|)}a_{<k}a_{k+1}a_{>k+1}\{a_k,b_j\}b_{\hat\jmath}$, while the term with $i=k+1$ of the sum for the interchanged arguments is $(-1)^{|a_k||a_{>k+1}|}a_{<k}a_{k+1}a_{>k+1}\{a_k,b_j\}b_{\hat\jmath}$; these differ by the factor $(-1)^{|a_k||a_{k+1}|}$, and the same holds for the pair of terms with $i=k+1$ and $i=k$.
For $R_m$ the computation is the same with $\{1,a_i\}$ in place of $\{a_i,b_j\}b_{\hat\jmath}$.

To verify (b), let $a_m=1$ and $a'=a_1\cdots a_{m-1}$.
In \eqref{Ddef} the term with $i=m$ equals $(-1)^{|a'|}a'\{1,1\}=\{1,1\}a'$ (recall that $\{1,1\}$ is odd), so $R_m(a_1,\dots,a_{m-1},1)=R_{m-1}(a_1,\dots,a_{m-1})+\{1,1\}a'-\{1,1\}a'$.
In \eqref{brdef} the terms with $i=m$ contribute $a'\sum_{j}(-1)^{|b_j||b_{<j}|}\{1,b_j\}b_{\hat\jmath}$, and since $(-1)^{|b_j||b_{<j}|}\{1,b_j\}b_{\hat\jmath}=(-1)^{|b_{<j}|}b_{<j}\{1,b_j\}b_{>j}$, this equals $a'\bigl(\D(b)+(n-1)\D(1)b\bigr)$ by \eqref{Ddef}; the terms with $i<m$ give the double sum for $a'$.
The remaining terms of \eqref{brdef} for $a$ and for $a'$ differ by $-a'\D(b)-(n-1)a'\D(1)b$, which cancels the extra contribution of the double sum.
The verification for $b_n=1$ is similar: writing $b'=b_1\cdots b_{n-1}$ and using $\{a_i,1\}=-(-1)^{\rp{a_i}}\{1,a_i\}$ (identity \eqref{id:skew} in $\Lh$), $a_{<i}\{1,a_i\}a_{>i}=(-1)^{\rp{a_i}|a_{>i}|}a_{\hat\imath}\{1,a_i\}$ and \eqref{Ddef}, one finds that the terms with $j=n$ of the double sum contribute $(-1)^{|a|}\bigl(\D(a)+(m-1)\D(1)a\bigr)b'$, which cancels the difference $(-1)^{\rp a}\D(a)b'-(m-1)\,a\D(1)b'$ between the remaining terms of \eqref{brdef} for $b$ and for $b'$.

Now define $\D$ on the basis $U$ of $\Sm$ and $\{\,,\}^*$ on $U\times U$ by \eqref{Ddef} and \eqref{brdef}, using for $u=e_{i_1}^{k_1}\cdots e_{i_r}^{k_r}\in U$ the factorization into $\deg_e(u)$ factors from $M\setminus\{1\}$ taken in nondecreasing order, and for $u=1$ the factorization with $m=1$, $a_1=1$; extend by linearity.
Then \eqref{Ddef} and \eqref{brdef} hold for arbitrary factorizations: both sides are multilinear in the factors, so it suffices to consider factors in the basis $M$ of $\Lh$; by (a) the factors can be rearranged in nondecreasing order, both sides changing by the same sign; if some odd basis element occurs twice, the left-hand side is zero and so is the right-hand side by (a), because it changes sign under the interchange of two equal odd arguments; by (b) the factors equal to $1$ can be removed, and what remains is the factorization used in the definition.

Statements (i), (ii) and (iv) follow directly from \eqref{Ddef} and \eqref{brdef} with $m=n=1$, respectively with $m=1$, $a_1=1$ (using the identity $(-1)^{|b_j||b_{<j}|}\{1,b_j\}b_{\hat\jmath}=(-1)^{|b_{<j}|}b_{<j}\{1,b_j\}b_{>j}$ established above and $(-1)^{\rp 1}=-1$).
For (iii), apply \eqref{Ddef} to the factorization $bc=b_1\cdots b_nc_1\cdots c_p$: the terms corresponding to the factors of $b$ give $(\D(b)+(n-1)\D(1)b)c$, the terms corresponding to the factors of $c$ give $(-1)^{|b|}b(\D(c)+(p-1)\D(1)c)=(-1)^{|b|}b\D(c)+(p-1)\D(1)bc$, and the last term is $-(n+p-1)\D(1)bc$.
\end{proof}

We now verify the identities \eqref{id:skew} and \eqref{id:leibniz} for $\{\,,\}^*$.
In the proofs we write $\{a,b\}_0$ for the double sum in \eqref{brdef}, so that
\begin{equation}\label{brsplit}
\{a,b\}^*=\{a,b\}_0-(m-1)\,a\D(b)+(n-1)(-1)^{\rp a}\D(a)b-(m-1)(n-1)\,a\D(1)b .
\end{equation}

\begin{lemma}\label{lem:skew}
For all $a,b\in\Sm$ we have $\{a,b\}^*=-(-1)^{\rp a\rp b}\{b,a\}^*$.
\end{lemma}

\begin{proof}
Let $a=a_1\cdots a_m$, $b=b_1\cdots b_n$ with $a_i,b_j\in\Lh$.
Consider the term of $\{b,a\}_0$ corresponding to the pair $(j,i)$:
$$
(-1)^{|b_j||b_{>j}|+|a_i||a_{<i}|}\,b_{\hat\jmath}\{b_j,a_i\}a_{\hat\imath}.
$$
By \eqref{id:skew} in $\Lh$, $\{b_j,a_i\}=-(-1)^{\rp{a_i}\rp{b_j}}\{a_i,b_j\}$, and by \eqref{id:comm}
$$
b_{\hat\jmath}\{a_i,b_j\}a_{\hat\imath}=(-1)^{|a_{\hat\imath}|(|b_{\hat\jmath}|+\rp{a_i}+|b_j|)+|b_{\hat\jmath}|(\rp{a_i}+|b_j|)}\,a_{\hat\imath}\{a_i,b_j\}b_{\hat\jmath},
$$
since $|\{a_i,b_j\}|=\rp{a_i}+|b_j|$.
Using $\rp a=|a_{\hat\imath}|+\rp{a_i}$, $\rp b=|b_{\hat\jmath}|+\rp{b_j}$, $|b_{\hat\jmath}|=|b|+|b_j|$, $|a_{\hat\imath}|=|a|+|a_i|$ and $|a_i|+\rp{a_i}=|b_j|+\rp{b_j}=1$, a direct comparison of the exponents shows that this term equals
$$
-(-1)^{\rp a\rp b}\,(-1)^{|a_i||a_{>i}|+|b_j||b_{<j}|}\,a_{\hat\imath}\{a_i,b_j\}b_{\hat\jmath},
$$
i.e. $-(-1)^{\rp a\rp b}$ times the $(i,j)$-term of $\{a,b\}_0$.
Hence $\{b,a\}_0=-(-1)^{\rp a\rp b}\{a,b\}_0$.

For the remaining terms of \eqref{brsplit} we use $|\D(a)|=\rp a$, $|\D(1)|=1$ and \eqref{id:comm}:
$$
b\D(a)=(-1)^{|b|\rp a}\D(a)b,\qquad \D(b)a=(-1)^{\rp b|a|}a\D(b),\qquad b\D(1)a=-(-1)^{\rp a\rp b}a\D(1)b .
$$
Therefore
\begin{multline*}
-(n-1)\,b\D(a)+(m-1)(-1)^{\rp b}\D(b)a-(n-1)(m-1)\,b\D(1)a\\
=-(-1)^{\rp a\rp b}\Bigl(-(m-1)\,a\D(b)+(n-1)(-1)^{\rp a}\D(a)b-(m-1)(n-1)\,a\D(1)b\Bigr),
\end{multline*}
because $(-1)^{\rp b(|a|+1)}=(-1)^{\rp a\rp b}$ and $(-1)^{|b|\rp a}=(-1)^{\rp a\rp b+\rp a}$.
Adding the two computations we obtain the lemma.
\end{proof}

\begin{lemma}\label{lem:leibniz}
For all $a,b,c\in\Sm$ we have
$$
\{a,bc\}^*=\{a,b\}^*c+(-1)^{\rp a|b|}b\{a,c\}^*+(-1)^{\rp a}\D(a)bc .
$$
\end{lemma}

\begin{proof}
Let $a=a_1\cdots a_m$, $b=b_1\cdots b_n$, $c=c_1\cdots c_p$ with $a_i,b_j,c_k\in\Lh$, and apply \eqref{brdef} to the factorization $bc=b_1\cdots b_nc_1\cdots c_p$.
The terms of $\{a,bc\}_0$ corresponding to the factors of $b$ give $\{a,b\}_0c$.
A term corresponding to the factor $c_k$ equals
$$
(-1)^{|a_i||a_{>i}|+|c_k|(|b|+|c_{<k}|)}\,a_{\hat\imath}\{a_i,c_k\}\,b\,c_{\hat k}
=(-1)^{\rp a|b|}\,(-1)^{|a_i||a_{>i}|+|c_k||c_{<k}|}\,b\,a_{\hat\imath}\{a_i,c_k\}c_{\hat k},
$$
because $|a_{\hat\imath}\{a_i,c_k\}|=\rp a+|c_k|$; hence these terms give $(-1)^{\rp a|b|}b\{a,c\}_0$.
Thus
$$
\{a,bc\}_0=\{a,b\}_0c+(-1)^{\rp a|b|}b\{a,c\}_0 .
$$
It remains to compare the terms involving $\D$.
By \eqref{brsplit}, on the left-hand side they are
$$
-(m-1)\,a\D(bc)+(n+p-1)(-1)^{\rp a}\D(a)bc-(m-1)(n+p-1)\,a\D(1)bc,
$$
and on the right-hand side they are
\begin{multline*}
\bigl(-(m-1)\,a\D(b)+(n-1)(-1)^{\rp a}\D(a)b-(m-1)(n-1)\,a\D(1)b\bigr)c\\
+(-1)^{\rp a|b|}b\bigl(-(m-1)\,a\D(c)+(p-1)(-1)^{\rp a}\D(a)c-(m-1)(p-1)\,a\D(1)c\bigr)+(-1)^{\rp a}\D(a)bc .
\end{multline*}
Using \eqref{id:comm} together with $|\D(a)|=\rp a$ and $|\D(1)|=1$ we have
$$
(-1)^{\rp a|b|}ba\D(c)=(-1)^{|b|}ab\D(c),\quad (-1)^{\rp a|b|}b\D(a)c=\D(a)bc,\quad (-1)^{\rp a|b|}ba\D(1)c=a\D(1)bc,
$$
so the right-hand side equals
$$
-(m-1)\,a\bigl(\D(b)c+(-1)^{|b|}b\D(c)\bigr)+(n+p-1)(-1)^{\rp a}\D(a)bc-(m-1)(n+p-2)\,a\D(1)bc,
$$
which coincides with the left-hand side by Lemma~\ref{lem:welldef}(iii).
\end{proof}

To prove the Jacobi identity we use the following general formula for the Jacobiator.

\begin{lemma}\label{jacobiator}
Let $(V,\cdot,\{\,,\})$ be a superalgebra with unit satisfying \eqref{id:deg}, \eqref{id:comm}, \eqref{id:skew} and \eqref{id:leibniz}, and let
$$
J(a,b,c)=\{a,\{b,c\}\}-\{\{a,b\},c\}-(-1)^{\rp a\rp b}\{b,\{a,c\}\}
$$
for $a,b,c\in V$.
Then
\begin{enumerate}
\item[(i)] $J(a,b,c)=-(-1)^{\rp a\rp b}J(b,a,c)=-(-1)^{\rp b\rp c}J(a,c,b)$;
\item[(ii)] $J(a,b,xy)=J(a,b,x)\,y+(-1)^{(|a|+|b|)|x|}\,x\,J(a,b,y)-J(a,b,1)\,xy$ for all $a,b,x,y\in V$.
\end{enumerate}
\end{lemma}

\begin{proof}
(i) By \eqref{id:skew}, $\{\{b,a\},c\}=-(-1)^{\rp a\rp b}\{\{a,b\},c\}$, whence
$$
-(-1)^{\rp a\rp b}J(b,a,c)=-(-1)^{\rp a\rp b}\{b,\{a,c\}\}-\{\{a,b\},c\}+\{a,\{b,c\}\}=J(a,b,c).
$$
Similarly, $\{a,\{c,b\}\}=-(-1)^{\rp b\rp c}\{a,\{b,c\}\}$, $\{\{a,c\},b\}=-(-1)^{\rp b(\rp a+\rp c)}\{b,\{a,c\}\}$ and $\{c,\{a,b\}\}=-(-1)^{\rp c(\rp a+\rp b)}\{\{a,b\},c\}$, since $\rp{\{a,c\}}=\rp a+\rp c$ and $\rp{\{a,b\}}=\rp a+\rp b$; substituting these into $J(a,c,b)$ gives $J(a,c,b)=-(-1)^{\rp b\rp c}J(a,b,c)$.

(ii) We expand the three terms of $J(a,b,xy)$ by \eqref{id:leibniz}, using $|\{b,x\}|=\rp b+|x|$, $|\D(b)|=\rp b$, $\rp{\{a,b\}}=\rp a+\rp b$ and $(-1)^{(\rp a+\rp b)|x|}=(-1)^{(|a|+|b|)|x|}$:
\begin{align*}
\{a,\{b,xy\}\}={}&\{a,\{b,x\}\}y+(-1)^{\rp a(|x|+\rp b)}\{b,x\}\{a,y\}+(-1)^{\rp a}\D(a)\{b,x\}y\\
&+(-1)^{\rp b|x|}\Bigl(\{a,x\}\{b,y\}+(-1)^{\rp a|x|}x\{a,\{b,y\}\}+(-1)^{\rp a}\D(a)x\{b,y\}\Bigr)\\
&+(-1)^{\rp b}\Bigl(\{a,\D(b)\}xy+(-1)^{\rp a\rp b}\D(b)\{a,xy\}+(-1)^{\rp a}\D(a)\D(b)xy\Bigr),\\
\{\{a,b\},xy\}={}&\{\{a,b\},x\}y+(-1)^{(|a|+|b|)|x|}x\{\{a,b\},y\}+(-1)^{\rp a+\rp b}\D(\{a,b\})xy,\\
\{b,\{a,xy\}\}={}&\{b,\{a,x\}\}y+(-1)^{\rp b(|x|+\rp a)}\{a,x\}\{b,y\}+(-1)^{\rp b}\D(b)\{a,x\}y\\
&+(-1)^{\rp a|x|}\Bigl(\{b,x\}\{a,y\}+(-1)^{\rp b|x|}x\{b,\{a,y\}\}+(-1)^{\rp b}\D(b)x\{a,y\}\Bigr)\\
&+(-1)^{\rp a}\Bigl(\{b,\D(a)\}xy+(-1)^{\rp a\rp b}\D(a)\{b,xy\}+(-1)^{\rp b}\D(b)\D(a)xy\Bigr),
\end{align*}
where, in addition, $\{a,xy\}$ and $\{b,xy\}$ are to be expanded by \eqref{id:leibniz}.
We substitute these expansions into $J(a,b,xy)=\{a,\{b,xy\}\}-\{\{a,b\},xy\}-(-1)^{\rp a\rp b}\{b,\{a,xy\}\}$.
The terms $\{a,\{b,x\}\}y$, $-\{\{a,b\},x\}y$, $-(-1)^{\rp a\rp b}\{b,\{a,x\}\}y$ add up to $J(a,b,x)y$, and the terms containing $x\{a,\{b,y\}\}$, $x\{\{a,b\},y\}$, $x\{b,\{a,y\}\}$ add up to $(-1)^{(|a|+|b|)|x|}xJ(a,b,y)$.
By \eqref{id:skew} we have $\{b,1\}=-(-1)^{\rp b}\D(b)$, $\{a,1\}=-(-1)^{\rp a}\D(a)$ and $\{\{a,b\},1\}=-(-1)^{\rp a+\rp b}\D(\{a,b\})$, so that
$$
J(a,b,1)=-(-1)^{\rp b}\{a,\D(b)\}+(-1)^{\rp a+\rp b}\D(\{a,b\})+(-1)^{\rp a\rp b+\rp a}\{b,\D(a)\};
$$
hence the terms $(-1)^{\rp b}\{a,\D(b)\}xy$, $-(-1)^{\rp a+\rp b}\D(\{a,b\})xy$, $-(-1)^{\rp a\rp b+\rp a}\{b,\D(a)\}xy$ add up to $-J(a,b,1)xy$.
The remaining sixteen terms cancel in pairs: after the expansion of $\{a,xy\}$ and $\{b,xy\}$ they add up to
\begin{gather*}
\bigl((-1)^{\rp a(|x|+\rp b)}-(-1)^{\rp a\rp b+\rp a|x|}\bigr)\{b,x\}\{a,y\}
+\bigl((-1)^{\rp b|x|}-(-1)^{\rp a\rp b+\rp b(|x|+\rp a)}\bigr)\{a,x\}\{b,y\}\\
+\bigl((-1)^{\rp a}-(-1)^{\rp a}\bigr)\D(a)\{b,x\}y
+\bigl((-1)^{\rp b|x|+\rp a}-(-1)^{\rp a+\rp b|x|}\bigr)\D(a)x\{b,y\}\\
+\bigl((-1)^{\rp b+\rp a\rp b}-(-1)^{\rp a\rp b+\rp b}\bigr)\D(b)\{a,x\}y\\
+\bigl((-1)^{\rp b+\rp a\rp b+\rp a|x|}-(-1)^{\rp a\rp b+\rp a|x|+\rp b}\bigr)\D(b)x\{a,y\}\\
+\bigl((-1)^{\rp b+\rp a}-(-1)^{\rp a+\rp b}\bigr)\D(a)\D(b)xy
+\bigl((-1)^{\rp b+\rp a\rp b+\rp a}-(-1)^{\rp a\rp b+\rp a+\rp b}\bigr)\D(b)\D(a)xy,
\end{gather*}
and each of the eight coefficients is zero.
\end{proof}

\begin{proposition}\label{prop:jacobi}
The bracket $\{\,,\}^*$ satisfies the odd Jacobi identity \eqref{id:jacobi} on $\Sm$.
\end{proposition}

\begin{proof}
By Lemmas~\ref{lem:welldef}, \ref{lem:skew} and \ref{lem:leibniz}, the superalgebra $(\Sm,\cdot,\{\,,\}^*)$ satisfies the hypotheses of Lemma~\ref{jacobiator}; let $J$ be its Jacobiator.
We have to show that $J(a,b,c)=0$ for all $a,b,c\in\Sm$.

First, for fixed $a,b\in\Sm$, if $J(a,b,u)=0$ for all $u\in\Lh$, then $J(a,b,c)=0$ for all $c\in\Sm$.
Indeed, $\Sm$ is spanned by the products $c=c_1\cdots c_m$ of elements of $\Lh$, and by Lemma~\ref{jacobiator}(ii),
\begin{multline*}
J(a,b,c_1\cdots c_m)=J(a,b,c_1\cdots c_{m-1})\,c_m\\
+(-1)^{(|a|+|b|)|c_1\cdots c_{m-1}|}\,c_1\cdots c_{m-1}\,J(a,b,c_m)-J(a,b,1)\,c_1\cdots c_m,
\end{multline*}
so the claim follows by induction on $m$, since $1,c_m\in\Lh$.

Now let $u,v,w\in\Lh$.
By Lemma~\ref{lem:welldef}(i), $J(u,v,w)$ is the Jacobiator of the odd Lie superalgebra $\Lh$, hence $J(u,v,w)=0$.
By the previous paragraph, $J(u,v,c)=0$ for all $c\in\Sm$.
By Lemma~\ref{jacobiator}(i), $J(u,c,v)=-(-1)^{\rp c\rp v}J(u,v,c)=0$ for $u,v\in\Lh$, $c\in\Sm$; applying the previous paragraph once more we get $J(u,b,c)=0$ for all $u\in\Lh$ and $b,c\in\Sm$.
Finally, by Lemma~\ref{jacobiator}(i), $J(a,b,u)=-(-1)^{\rp b\rp u}J(a,u,b)=(-1)^{\rp b\rp u+\rp a\rp u}J(u,a,b)=0$ for $a,b\in\Sm$, $u\in\Lh$, and the first paragraph gives $J(a,b,c)=0$ for all $a,b,c\in\Sm$.
\end{proof}

Summarizing, we obtain

\begin{theorem}\label{thm:model}
$(\Sm,\cdot,\{\,,\}^*)$ is a unital generalized Gerstenhaber algebra, whose bracket restricted to $\Lh$ is the bracket of the free odd Lie superalgebra $\LL$ and whose map $\D$ is given by \eqref{Ddef}.
It is generated by $X$ as a unital generalized Gerstenhaber algebra.
\end{theorem}

\begin{proof}
The identities \eqref{id:deg} and \eqref{id:comm} hold in $\Sm$ by construction and by Lemma~\ref{lem:welldef}(iv); \eqref{id:skew}, \eqref{id:leibniz} and \eqref{id:jacobi} are Lemma~\ref{lem:skew}, Lemma~\ref{lem:leibniz} (note that $\D(a)=\{1,a\}^*$ by Lemma~\ref{lem:welldef}(ii)) and Proposition~\ref{prop:jacobi}.
The subalgebra of $\Sm$ generated by $X$ contains $\Lh$, because $\LL$ is generated by $X$ as an odd Lie superalgebra and $\{\,,\}^*$ coincides with the bracket of $\LL$ on $\Lh$; and $\Lh$ generates $\Sm$ under the product.
\end{proof}

\section{Main results}

\begin{theorem}\label{basis}
The set $U$ is a basis of the free unital generalized Gerstenhaber algebra $\G$ generated by $X$.
Moreover, the homomorphism $\varphi\colon\G\to\Sm$ which is the identity on $X$ is an isomorphism of generalized Gerstenhaber algebras.
\end{theorem}

\begin{proof}
By Theorem~\ref{thm:model} and the universal property of $\G$, the identity map on $X$ extends to a homomorphism $\varphi\colon\G\to\Sm$ of unital generalized Gerstenhaber algebras, which is surjective because $\Sm$ is generated by $X$.
Every element $e\in M$ is a word in the generators with respect to the bracket, and $\varphi$ is compatible with the brackets; hence $\varphi(e)=e$ for $e\in M$, and consequently $\varphi(u)=u$ for every $u\in U$.
Since $U$ is a basis of $\Sm$, the set $U\subset\G$ is linearly independent, and by Lemma~\ref{vsp} it spans $\G$.
Thus $U$ is a basis of $\G$, and $\varphi$, being a surjection mapping a basis bijectively onto a basis, is an isomorphism.
\end{proof}

\begin{corollary}\label{cor:11}
The element $\D(1)=\{1,1\}$ is nonzero in $\G$, and $\D$ is not a derivation of $(\G,\cdot)$.
More precisely, the elements $\{1,1\}$, $\{1,1\}x_i$ and $\{1,x_i\}$, $i\ge1$, are linearly independent in $\G$.
\end{corollary}

\begin{proof}
The words $\{1,1\}$ (a square of the even good word $1$), $\{x_i,1\}$ (a good word) and the monomials $\{1,1\}x_i$ belong to $U$; note that $\{1,x_i\}=-(-1)^{\rp{x_i}}\{x_i,1\}$.
The second assertion follows from Remark~\ref{rem:D}(a).
\end{proof}

Following~\cite{K17}, for $n\ge1$ consider the free unital generalized Gerstenhaber algebra $\G$ generated by $X=\{1,x_1,\dots,x_n\}$ (with arbitrary parities of $x_1,\dots,x_n$) and denote by $\mathrm{GenGer}(n)$ the subspace of $\G$ spanned by the elements of the basis $U$ which are multilinear in $1,x_1,\dots,x_n$, that is, in which each of the letters $1,x_1,\dots,x_n$ occurs exactly once.

\begin{theorem}\label{dimen}
$\dim\mathrm{GenGer}(n)=n\cdot n!$.
\end{theorem}

\begin{proof}
An element of $U$ is a product $e_{i_1}\cdots e_{i_r}$ of elements of $M\setminus\{1\}$, and it is multilinear in $1,x_1,\dots,x_n$ if and only if the factors are good words (squares $\{v,v\}$ contain every letter of $v$ twice) whose sets of letters are pairwise disjoint and multilinear, with union $X$.
Thus the multilinear elements of $U$ are in bijection with the pairs consisting of a partition $\pi$ of the set $X$ into blocks none of which equals $\{1\}$ (the factor $1$ is excluded in $U$) and a family of multilinear good words $w_B$, $B\in\pi$, where the set of letters of $w_B$ is $B$.

The multilinear good words with a given set $B$ of $k$ letters form, by Lemma~\ref{lem:M}, a basis of the multilinear component of the free odd Lie superalgebra generated by $B$; since good words are defined purely combinatorially, their number does not depend on the parities of the letters, and it is equal to the dimension $(k-1)!$ of the multilinear component of the free Lie algebra of rank $k$ (see~\cite{Reu93}).
Consequently,
$$
\dim\mathrm{GenGer}(n)=\sum_{\pi}\ \prod_{B\in\pi}(|B|-1)!,
$$
where $\pi$ runs over the partitions of $X$ without the block $\{1\}$.
For every finite set $Y$ we have $\sum_{\pi}\prod_{B\in\pi}(|B|-1)!=|Y|!$, the sum being over all partitions of $Y$: a partition of $Y$ together with a cyclic order on each block is the same as a permutation of $Y$ (the blocks being its cycles).
Subtracting the partitions of $X$ containing the block $\{1\}$, which correspond to the partitions of $\{x_1,\dots,x_n\}$, we obtain $\dim\mathrm{GenGer}(n)=(n+1)!-n!=n\cdot n!$.
\end{proof}

\begin{remark}
The same argument, applied to the free odd Lie superalgebra $\LL_0$ generated by $\{x_1,x_2,\dots\}$ with the bracket extended to $\F1\oplus\LL_0$ by $\{1,\cdot\}=\{\cdot,1\}=0$, gives $\D=0$ in \eqref{Ddef} and turns \eqref{brdef} into the unique extension of the bracket of $\LL_0$ to a biderivation of the free supercommutative superalgebra generated by $\LL_0$.
Lemmas~\ref{lem:welldef}--\ref{jacobiator} and Proposition~\ref{prop:jacobi} remain valid verbatim, and the proof of Lemma~\ref{vsp} goes through with $\{1,u\}=0$.
This recovers the well-known description of the free unital Gerstenhaber algebra generated by $x_1,x_2,\dots$: its basis is formed by the elements of $U$ in which the letter $1$ does not occur.
\end{remark}

\begin{thebibliography}{99}

\bibitem{BMPZ}
Bahturin Yu., Mikhalev A., Petrogradsky V., Zaicev M.,
Infinite-dimensional Lie superalgebras,
De Gruyter Expositions in Mathematics, 7. Walter de Gruyter \& Co., Berlin, 1992.

\bibitem{BV81}
Batalin I., Vilkovisky G.,
Gauge algebra and quantization,
Physics Letters B, 102 (1981), 1, 27--31.

\bibitem{p2}
Batalin I., Bering K.,
Odd scalar curvature in anti-Poisson geometry,
Physics Letters B, 663 (2008), 1--2, 132--135.

\bibitem{BLLM}
Bell J., Launois S., Le\'on S\'anchez O., Moosa R.,
Poisson algebras via model theory and differential algebraic geometry,
Journal of the European Mathematical Society, 19 (2017), 7, 2019--2049.

\bibitem{BKLM}
Bokut L., Kang S.-J., Lee K.-H., Malcolmson P.,
Gr\"obner--Shirshov bases for Lie superalgebras and their universal enveloping algebras,
Journal of Algebra, 217 (1999), 2, 461--495.

\bibitem{B06}
Bratchikov A.,
Subalgebras of the Gerstenhaber algebra of differential operators,
Journal of Algebra and its Applications, 16 (2017), 1, 1750005, 6 pp.

\bibitem{p1}
Bruce A., Grabowski J.,
Odd connections on supermanifolds: existence and relation with affine connections,
Journal of Physics A, 53 (2020), 45, 455203, 24 pp.

\bibitem{CK07}
Cantarini N., Kac V.,
Classification of linearly compact simple Jordan and generalized Poisson superalgebras,
Journal of Algebra, 313 (2007), 1, 100--124.

\bibitem{CK10}
Cantarini N., Kac V.,
Classification of linearly compact simple rigid superalgebras,
International Mathematics Research Notices, (2010), 17, 3341--3393.

\bibitem{chi}
Chibrikov E.,
A right normed basis for free Lie algebras and Lyndon--Shirshov words,
Journal of Algebra, 302 (2006), 2, 593--612.

\bibitem{Ger63}
Gerstenhaber M.,
The cohomology structure of an associative ring,
Annals of Mathematics, 78 (1963), 2, 267--288.

\bibitem{GK04}
Ginzburg V., Kaledin D.,
Poisson deformations of symplectic quotient singularities,
Advances in Mathematics, 186 (2004), 1, 1--57.

\bibitem{Hall50}
Hall M.,
A basis for free Lie rings and higher commutators in free groups,
Proceedings of the American Mathematical Society, 1 (1950), 575--581.

\bibitem{Kac77}
Kac V.,
Lie superalgebras,
Advances in Mathematics, 26 (1977), 1, 8--96.

\bibitem{Kan90}
Kantor I.,
Connection between Poisson brackets and Jordan and Lie superalgebras,
Lie theory, differential equations and representation theory (Montreal, 1989), 213--225, Univ. Montr\'eal, Montreal, 1990.

\bibitem{Kan92}
Kantor I.,
Jordan and Lie superalgebras determined by a Poisson algebra,
American Mathematical Society Translations, Series 2, 151 (1992), 55--80.

\bibitem{K10}
Kaygorodov I.,
Generalized Kantor double,
Vestnik Samarskogo Gosudarstvennogo Universiteta, 78 (2010), 4, 42--50.

\bibitem{K17}
Kaygorodov I.,
Algebras of Jordan brackets and generalized Poisson algebras,
Linear and Multilinear Algebra, 65 (2017), 6, 1142--1157.

\bibitem{KSU}
Kaygorodov I., Shestakov I., Umirbaev U.,
Free generic Poisson fields and algebras,
Communications in Algebra, 46 (2018), 4, 1799--1812.

\bibitem{KM92}
King D., McCrimmon K.,
The Kantor construction of Jordan superalgebras,
Communications in Algebra, 20 (1992), 1, 109--126.

\bibitem{Kon03}
Kontsevich M.,
Deformation quantization of Poisson manifolds,
Letters in Mathematical Physics, 66 (2003), 3, 157--216.

\bibitem{K96}
Kosmann-Schwarzbach Y.,
From Poisson algebras to Gerstenhaber algebras,
Annales de l'Institut Fourier (Grenoble), 46 (1996), 5, 1243--1274.

\bibitem{L1}
Li L.-C.,
Classical $r$-matrices and compatible Poisson structures for Lax equations on Poisson algebras,
Communications in Mathematical Physics, 203 (1999), 3, 573--592.

\bibitem{L21}
Lopes S., Solotar A.,
Lie structure on the Hochschild cohomology of a family of subalgebras of the Weyl algebra,
Journal of Noncommutative Geometry, 15 (2021), 4, 1373--1407.

\bibitem{Reu93}
Reutenauer C.,
Free Lie algebras,
London Mathematical Society Monographs. New Series, 7. Oxford University Press, New York, 1993.

\bibitem{She93}
Shestakov I.,
Quantization of Poisson superalgebras and speciality of Jordan Poisson superalgebras,
Algebra and Logic, 32 (1993), 5, 309--317.

\bibitem{shirshov}
Shirshov A. I.,
Selected works of A. I. Shirshov,
Translated from the Russian by M. Bremner and M. V. Kotchetov. Edited by L. A. Bokut, V. Latyshev, I. Shestakov and E. Zelmanov. Contemporary Mathematicians. Birkh\"auser Verlag, Basel, 2009.

\bibitem{S86}
Shtern A.,
Free Lie superalgebras,
Siberian Mathematical Journal, 27 (1986), 1, 136--140.

\bibitem{ber08}
Van den Bergh M.,
Double Poisson algebras,
Transactions of the American Mathematical Society, 360 (2008), 11, 5711--5769.

\bibitem{V16}
Volkov Yu.,
BV differential on Hochschild cohomology of Frobenius algebras,
Journal of Pure and Applied Algebra, 220 (2016), 10, 3384--3402.

\bibitem{V19}
Volkov Yu.,
Gerstenhaber bracket on the Hochschild cohomology via an arbitrary resolution,
Proceedings of the Edinburgh Mathematical Society (2), 62 (2019), 3, 817--836.

\bibitem{W21}
Wang Z.,
Gerstenhaber algebra and Deligne's conjecture on the Tate--Hochschild cohomology,
Transactions of the American Mathematical Society, 374 (2021), 7, 4537--4577.

\bibitem{p3}
Xu P.,
Gerstenhaber algebras and BV-algebras in Poisson geometry,
Communications in Mathematical Physics, 200 (1999), 3, 545--560.

\end{thebibliography}

\end{document}
